(1) 令稳态解 \(x(t)=\Re[\tilde X e^{i\omega t}]\)。代入方程得到复振幅
\[
\tilde X=\frac{F_0}{k-m\omega^2+i b\omega}.
\]
写成 \(x=X\cos(\omega t-\delta)\)，则
\[
X=\frac{F_0}{\sqrt{(k-m\omega^2)^2+b^2\omega^2}},\qquad
\delta=\operatorname{atan2}(b\omega,k-m\omega^2)\in(0,\pi).
\]
这种相位取值能正确描述穿过共振时相位由零趋向 \(\pi\)。

(2) 平均输入功率由 \(\overline P=\langle F_0\cos\omega t\,\dot x\rangle\) 得
\[
\overline P(\omega)=\frac{F_0^2 b\omega^2}{2[(k-m\omega^2)^2+b^2\omega^2]}.
\]
也可用耗散功率 \(b\langle\dot x^2\rangle=b\omega^2X^2/2\) 独立计算。由于分母不小于 \(b^2\omega^2\)，\(\overline P\le F_0^2/(2b)\)；等号恰在 \(k-m\omega^2=0\)，即 \(\omega=\omega_0\) 取得。

(3) 对振幅分母的平方 \(D(\omega)=(k-m\omega^2)^2+b^2\omega^2\) 求导，
\[
D'(\omega)=2\omega[2m(m\omega^2-k)+b^2].
\]
当 \(b^2<2km\)，非零的振幅峰频率
\[
\omega_X=\sqrt{\frac{k}{m}-\frac{b^2}{2m^2}}<\omega_0.
\]
若 \(b^2\ge2km\)，\(D\) 对 \(\omega>0\) 单调非降，振幅随频率下降（最大静态极限在 \(\omega\to0^+\)），但功率峰仍位于 \(\omega_0\)。两者物理机制不同：一个比较位移，一个比较能量输运。

(4) 半功率方程 \(\overline P(\omega)=F_0^2/(4b)\) 等价于
\[
(k-m\omega^2)^2=b^2\omega^2 .
\]
两条正频率解恰为
\[
\omega_-=\frac{\sqrt{b^2+4km}-b}{2m},\qquad
\omega_+=\frac{\sqrt{b^2+4km}+b}{2m}.
\]
故精确半高全宽 \(\Delta\omega=\omega_+-\omega_-=b/m\)，而不是仅在弱阻尼时才成立。弱阻尼 \(b\ll2\sqrt{km}\) 时 \(\omega_\pm\simeq\omega_0\pm b/(2m)\)；对应品质因子 \(Q_{\mathrm{osc}}=\omega_0/\Delta\omega=\sqrt{km}/b\)。
